% !TEX root = ../libretto.tex
% Sezione: Offertorio
\momento{Offertorio}
\begin{songs}{titoli}

% Fonte: LibrettoAccordi.pdf (Canti Coro), G6.
% Nell'elenco: Offertorio, Ordinario.
\beginsong{Accogli i nostri doni}
\beginverse
Accogli Signore i nostri doni
in questo misterioso incontro,
tra la nostra povertà e la tua grandezza.
\endverse
\beginverse
Noi ti offriamo le cose che tu stesso ci hai dato
e tu in cambio donaci, donaci te stesso.
\endverse
\endsong

% Fonte: LibrettoAccordi.pdf (Canti Coro), G11.
% Nell'elenco: Offertorio, Ordinario.
\beginsong{Antica eterna danza}[by={Gen Verde}]
\beginverse
Spighe d'oro al vento, antica eterna danza
per fare un solo pane spezzato sulla mensa.
Grappoli dei colli, profumo di letizia
per fare un solo vino, bevanda della grazia.
\endverse
\beginverse
Con il pane e il vino, Signore, ti doniamo
le nostre gioie pure, le attese e le paure,
frutti del lavoro e fede nel futuro,
la voglia di cambiare e di ricominciare.
\endverse
\beginverse
Dio della speranza, sorgente d'ogni dono,
accogli questa offerta che insieme ti portiamo.
Dio dell'universo, raccogli chi è disperso
e facci tutti Chiesa, una cosa in te.
\endverse
\endsong

% Fonte: LibrettoAccordi.pdf (Canti Coro), G15.
% Nell'elenco: Offertorio, Ordinario.
\beginsong{Benedici, o Signore}[by={Gen Rosso}]
\beginverse
Nebbia e freddo, giorni lunghi e amari mentre il seme muore.
Poi il prodigio antico e sempre nuovo del primo filo d'erba.
E nel vento dell'estate ondeggiano le spighe:
avremo ancora pane.
\endverse
\beginchorus
Benedici, o Signore,
questa offerta che portiamo a Te,
facci «uno» come il pane
che anche oggi hai dato a noi.
\endchorus
\beginverse
Nei filari dopo il lungo inverno fremono le viti.
La rugiada avvolge nel silenzio i primi tralci verdi.
Poi i colori dell'autunno coi grappoli maturi
avremo ancora vino.
\endverse
\endsong

% Fonte: LibrettoAccordi.pdf (Canti Coro), G10.
% Nell'elenco: Offertorio, Ordinario.
\beginsong{Come tanti fiumi}[by={Grasso, Selleri}]
\beginverse
Come tanti fiumi che nel mare si rispecchiano,
come tanti raggi che dal sole si diramano,
così noi, Signore Dio, dalle nostre mani
un solo pane ti offriamo, perché in te, in te
segno d'unità diventerà.
\endverse
\beginchorus
Fa' di noi, Signore, fiumi di speranza,
fa' che i nostri occhi vedano l'invisibile;
e nel volto dei fratelli e nelle nostre città
siamo semi di pace e d'unità.
\endchorus
\beginverse
Come dai sentieri una grande strada partirà,
come tante voci una parola nascerà,
così noi, Signore Dio, dalle nostre mani
un po' di vino ti offriamo, perché in te, in te
segno d'unità diventerà.
\endverse
\beginverse
Fa' di noi, Signore, offerta gradita a te,
facci vivere e sognare un mondo che
ama ogni uomo, e nella diversità
scopre te, che sei amore ed unità.
Scopre te, che sei amore ed unità.
\endverse
\endsong

% Fonte: LibrettoAccordi.pdf (Canti Coro), G9.
% Nell'elenco: Offertorio, Ordinario.
\beginsong{Cosa offrirti}[by={D. Branca, L. Ciancio}]
\beginverse
Cosa offrirti, o Dio, cosa posso darti,
eccomi son qui davanti a te.
Le gioie ed i dolori, gli affanni di ogni giorno,
tutto voglio vivere in te.
\endverse
\beginchorus
Accetta, mio re, questo poco che ho,
offro a Te la mia vita.
Gioia è per me far la tua volontà,
il mio unico bene sei solo Tu, solo Tu.
\endchorus
\beginverse
Vengo a te mio Dio, apro le mie braccia,
che la tua letizia riempirà.
Rinnova questo cuore perché ti sappia amare,
e nella tua pace io vivrò.
\endverse
\endsong

% Fonte: LibrettoAccordi.pdf (Canti Coro), G7.
% Nell'elenco: Offertorio, Ordinario.
\beginsong{Ecco quel che abbiamo}[by={Gen Verde}]
\beginverse
Ecco quel che abbiamo, nulla ci appartiene ormai.
Ecco i frutti della terra, che Tu moltiplicherai.
Ecco queste mani, puoi usarle se lo vuoi,
per dividere nel mondo il pane che Tu hai dato a noi.
\endverse
\beginverse
Solo una goccia hai messo fra le mani mie,
solo una goccia che Tu ora chiedi a me.
Una goccia che, in mano a Te, una pioggia diventerà
e la terra feconderà.
\endverse
\beginverse
Sulle strade il vento da lontano porterà
il profumo del frumento, che tutti avvolgerà.
E sarà l'amore che il raccolto spartirà,
e il miracolo del pane in terra si ripeterà.
\endverse
\beginverse
Le nostre gocce, pioggia tra le mani Tue,
saranno linfa di una nuova civiltà.
E la terra preparerà la festa del pane che
ogni uomo condividerà.
\endverse
\endsong

% Fonte: LibrettoAccordi.pdf (Canti Coro), G12.
% Nell'elenco: Offertorio, Ordinario.
\beginsong{Frutto della nostra terra}[by={F. Buttazzo}]
\beginverse
Frutto della nostra terra, del lavoro di ogni uomo:
pane della nostra vita, cibo della quotidianità.
Tu che lo prendevi un giorno, lo spezzavi per i tuoi,
oggi vieni in questo pane, cibo vero dell'umanità.
\endverse
\beginchorus
E sarò pane, e sarò vino, nella mia vita, nelle tue mani.
Ti accoglierò dentro di me farò di me un'offerta viva,
un sacrificio gradito a te.
\endchorus
\beginverse
Frutto della nostra terra, del lavoro di ogni uomo:
vino delle nostre vigne sulla mensa dei fratelli tuoi.
Tu che lo prendevi un giorno, lo bevevi con i tuoi,
oggi vieni in questo vino e ti doni per la vita mia.
\endverse
\endsong

% Fonte: testo fornito dall'utente (tolti accordi e stacchetti musicali).
% Nell'elenco: "I frutti che ti offriamo". Finale: ultimi due versi ripetuti.
\beginsong{I frutti che ti offriamo}[by={F. Volpintesta}]
\beginverse
Pane bianco della terra qui per te, Signore,
ed il frutto della vite ti doniamo lieti:
la tua forza benedica questi frutti che,
con l'impegno di ogni giorno,
noi offriamo a te.
\endverse
\beginverse
Ogni ora di lavoro, ogni affanno in più,
il cercare ogni giorno la tua volontà
sono istanti della vita, sono frutti che,
con l'impegno di ogni giorno,
noi offriamo a te.
\endverse
\beginverse
Tanti sbagli da bambino e la nostalgia,
una gioia, un dolore, quella povertà:
tu li accetti, o Signore, sono frutti che,
con l'impegno di ogni giorno,
noi offriamo a te.
\endverse
\beginverse*
Con l'impegno di ogni giorno
noi offriamo a te.
\endverse
\endsong

% Fonte: LibrettoAccordi.pdf (Canti Coro), G1.
% Nell'elenco: Offertorio, Ordinario (autore in elenco: Gen Verde).
\beginsong{Il seme del tuo campo}[by={Gen Verde; testo P. A. Sequeri}]
\beginverse
Per ogni volta che ci doni la parola di Luce, noi offriremo la pace.
Per ogni volta che ci nutre il tuo Pane di vita, noi sazieremo la fame.
Per ogni volta che ci allieta il tuo vino di gioia, noi guariremo ferite.
\endverse
\beginchorus
Offriamo a te, sinceramente, la vita, benediciamo la tua pace fra noi.
Saremo l'eco del tuo canto, il seme del tuo campo,
il lievito del tuo perdono, il lievito del tuo perdono.
\endchorus
\beginverse
Non ci separa dalla fede l'incertezza del cuore, quando ci parli, Signore.
Non ci separa dall'amore la potenza del male, quando rimani con noi.
Non ci separa dall'attesa del tuo giorno la morte, quando ci tieni per mano.
\endverse
\endsong

% Fonte: LibrettoAccordi.pdf (Canti Coro), G16; c'e' anche "6 - Offertorio/06 Quelli che amano Te.pdf".
% Nell'elenco: Offertorio, Ordinario. CD: Voce del mio canto.
\beginsong{Quelli che amano te}[by={Gen Rosso (CD Voce del mio canto)}]
\beginverse
Ecco il frumento che noi ti portiamo,
pane della madre terra
e del lavoro e sudore di tanti,
quelli che amano te.
\endverse
\beginverse
Portiamo il vino, frutto della vite,
sangue della madre terra
e del lavoro e la gioia di tanti,
quelli che amano te.
\endverse
\beginverse
Insieme ai doni veniamo noi stessi,
figli di tutta la terra,
per riscaldarci attorno alla mensa
che hai preparato per noi.
\endverse
\endsong

% Fonte: LibrettoAccordi.pdf (Canti Coro), G5.
% Nell'elenco: Offertorio, Ordinario. CD: Come fuoco vivo.
\beginsong{Segni del tuo amore}[by={Gen Rosso, Gen Verde (CD Come fuoco vivo)}]
\beginverse
Mille e mille grani nelle spighe d'oro
mandano fragranza e danno gioia al cuore,
quando, macinati, fanno un pane solo:
pane quotidiano, dono tuo, Signore.
\endverse
\beginchorus
Ecco il pane e il vino, segni del tuo amore.
Ecco questa offerta, accoglila, Signore:
tu di mille e mille cuori fai un cuore solo, un corpo solo in te
e il Figlio tuo verrà, vivrà ancora in mezzo a noi.
\endchorus
\beginverse
Mille grappoli maturi sotto il sole,
festa della terra, donano vigore,
quando da ogni perla stilla il vino nuovo:
vino della gioia, dono tuo, Signore.
\endverse
\endsong

% Fonte: LibrettoAccordi.pdf (Canti Coro), G14.
% Nell'elenco: Offertorio, Ordinario. CD: Giorno senza tramonto (2013).
\beginsong{Seme della terra}[by={M. Balduzzi (CD Giorno senza tramonto)}]
\beginverse
Seme chiuso nella terra, sotto il manto della neve,
dorme, e morendo forse non lo sa, spiga alta verso il sole diverrà.
\endverse
\beginverse
Uomo curvo sulla terra sotto il peso della vita,
grida, e morendo forse non lo sa, cielo alto come il sole diverrà.
\endverse
\beginverse
Uomo, seme della terra, sotto il manto delle stelle,
ama, e morendo forse non lo sa, vive un istante di immortalità.
\endverse
\endsong

% Fonte: LibrettoAccordi.pdf (Canti Coro), G3.
% Nell'elenco: Offertorio, Ordinario.
\beginsong{Servo per amore}[by={Gen Rosso}]
\beginverse
Una notte di sudore sulla barca in mezzo al mare
e mentre il cielo si imbianca già, tu guardi le tue reti vuote.
Ma la voce che ti chiama un altro mare ti mostrerà
e sulle rive di ogni cuore le tue reti getterai.
\endverse
\beginchorus
Offri la vita tua
come Maria ai piedi della croce
e sarai servo di ogni uomo, servo per amore, sacerdote dell'umanità.
\endchorus
\beginverse
Avanzavi nel silenzio fra le lacrime e speravi
che il seme sparso davanti a te cadesse sulla buona terra.
Ora il cuore tuo è in festa perché il grano biondeggia ormai,
è maturato sotto il sole, puoi riporlo nei granai.
\endverse
\endsong

% Fonte: LibrettoAccordi.pdf (Canti Coro), G13.
% Nell'elenco: Offertorio, Ordinario.
\beginsong{Signore della vita}[by={D. Ricci}]
\beginverse
Signore della vita, Padre. Benedetto sei tu, Dio dell'universo,
dalla tua bontà abbiamo ricevuto il pane,
dalle spighe che ci doni, dalle mani che ci hai dato.
E questo pane noi lo portiamo a te, Signore, accoglilo.
Diventerà il cibo della vita eterna, il corpo tuo, Signore.
\endverse
\beginverse
Signore della vita, Padre. Benedetto sei tu, Dio dell'universo,
dalla tua bontà abbiamo ricevuto il vino,
dalle viti che ci doni, dalle mani che ci hai dato.
E questo vino noi lo portiamo a te, Signore, accoglilo.
Diventerà la bevanda di salvezza, il sangue tuo, Signore.
\endverse
\beginverse*
Signore della vita, Padre. Benedetto sei tu, Dio dell'universo.
\endverse
\endsong

% Fonte: LibrettoAccordi.pdf (Canti Coro), G2. Omesso l'intro "Na, na, na...".
% Nell'elenco: Offertorio, Matrimonio.
\beginsong{Vita nuova nell'amore}[by={F. Buttazzo}]
\beginverse
Vita nuova germogliata in noi, come grano che ha piantato Dio,
come il pane per l'umanità: è l'amore suo, è l'amore in noi.
Vita nuova di felicità, vino buono per l'umanità,
con la forza dello Spirito sarà via di santità. \rep{2}
\endverse
\beginverse
È l'amore eterno di Gesù, per la Chiesa, amata sposa sua,
il mistero che si celebra, che risplende qui, che risplende in noi.
Il Signore ha chiamato noi all'amore nella fedeltà,
benedice i nostri passi qui nella sua comunità. \rep{2}
\endverse
\beginverse
È una casa che si costruirà sulla roccia della verità,
fondamento che non cede mai: il Signore in noi, il Signore in noi.
È famiglia che feconderà questo mondo di vitalità;
è promessa di felicità, vita piena data a noi. \rep{2}
\endverse
\endsong

% Fonte: LibrettoAccordi.pdf (Canti Coro), G4. Tolte le indicazioni Uomini/Donne.
% Nell'elenco: Offertorio, Ordinario.
\beginsong{Voi siete di Dio}[by={M. Balduzzi (CD Verbum Panis)}]
\beginverse
Tutte le stelle della notte,
le nebulose e le comete,
il sole su una ragnatela, è tutto vostro e voi siete di Dio.
\endverse
\beginverse
Tutte le rose della vita, il grano, i prati, i fili d'erba,
il mare, i fiumi, le montagne, è tutto vostro e voi siete di Dio.
\endverse
\beginverse
Tutte le musiche e le danze, i grattacieli, le astronavi,
i quadri, i libri, le culture, è tutto vostro e voi siete di Dio.
\endverse
\beginverse
Tutte le volte che perdono, quando sorrido, quando piango,
quando mi accorgo di chi sono, è tutto vostro e voi siete di Dio.
È tutto nostro e noi siamo di Dio.
\endverse
\endsong

\end{songs}
